\documentclass[tikz, border=0]{standalone}
\usepackage{amsmath}
\usepackage{luacode} % Pacote necessário para rodar Lua
\usetikzlibrary{calc,arrows.meta,patterns}

\definecolor{Red}{RGB}{ 159, 36, 48 }
\definecolor{Blue}{RGB}{ 103, 178, 255 }
\definecolor{LBlue}{RGB}{ 186, 220, 255 }
\definecolor{DBlue}{HTML}{132C53} 
\definecolor{gray}{HTML}{708090}
\definecolor{orange}{HTML}{ed9552}
\definecolor{orange2}{HTML}{e29d40}

% =====================================================================
% BLOCO LUA: Cálculo das equações e do Convex Hull (Monotone Chain)
% =====================================================================
\begin{luacode*}
-- Agora a função recebe os limites espaciais definidos no LaTeX
function draw_convex_hull(Lx, Ly)
    local pts = {}
    
    -- O limite matemático da equação é a = -2 e b = 2.
    -- Para isso bater exatamente na borda do seu desenho, a escala deve ser:
    local scale_x = Lx / 2.0
    local scale_y = Ly / 2.0
    
    -- Funções matemáticas originais
    local function calc_a(q, eps)
        return (2^(-q+1))*(q - 1) - 2*eps*(1 - 2^(-q+1))
    end
    
    local function calc_b(q, eps)
        return (2^(-q))*(q - 1)*(q - 2)*(1 - q/3) + 2*eps*(1 - (2^(-q))*(2 - q + q^2))
    end

    -- 1. Gerar os pontos das fronteiras com os fatores de escala independentes
    for q = 1, 40, 0.05 do
        table.insert(pts, {x = calc_a(q, 0) * scale_x, y = calc_b(q, 0) * scale_y})
        table.insert(pts, {x = calc_a(q, 1) * scale_x, y = calc_b(q, 1) * scale_y})
    end
    -- Inserir o ponto limite (canto superior esquerdo matemático)
    table.insert(pts, {x = -2 * scale_x, y = 2 * scale_y})

    -- 2. Algoritmo Monotone Chain (não muda)
    table.sort(pts, function(a, b)
        if math.abs(a.x - b.x) < 1e-7 then return a.y < b.y end
        return a.x < b.x
    end)

    local function cross(o, a, b)
        return (a.x - o.x) * (b.y - o.y) - (a.y - o.y) * (b.x - o.x)
    end

    local lower = {}
    for i = 1, #pts do
        while #lower >= 2 and cross(lower[#lower-1], lower[#lower], pts[i]) <= 0 do
            table.remove(lower)
        end
        table.insert(lower, pts[i])
    end

    local upper = {}
    for i = #pts, 1, -1 do
        while #upper >= 2 and cross(upper[#upper-1], upper[#upper], pts[i]) <= 0 do
            table.remove(upper)
        end
        table.insert(upper, pts[i])
    end

    table.remove(upper)
    table.remove(lower)
    
    local hull = {}
    for i = 1, #lower do table.insert(hull, lower[i]) end
    for i = 1, #upper do table.insert(hull, upper[i]) end

    -- 3. Gerar o código TikZ
    local tikz_path = "\\fill[orange] "
    for i = 1, #hull do
        tikz_path = tikz_path .. string.format("(%.4f, %.4f) -- ", hull[i].x, hull[i].y)
    end
    tikz_path = tikz_path .. "cycle;"
    
    -- Imprimir o comando no documento LaTeX
    tex.sprint(tikz_path)
end
\end{luacode*}
% =====================================================================
%\draw[pattern=north east lines, pattern color=blue] (0,0) rectangle (2,1);
\begin{document}

\begin{tikzpicture}[x=1em, y=1em,
        arrow/.style={-{Stealth[length=2mm, width=2mm]}, line width=1pt, black}
]

        \def\Lx{22}
        \def\Ly{12}

        \useasboundingbox (-25,-14.5) rectangle (23,12.5);

        \coordinate (center) at (0,0);

        \coordinate (topC) at (0,\Ly);
        \coordinate (topR) at (\Lx,\Ly);
        \coordinate (topL) at (-\Lx,\Ly);
        \coordinate (topRm) at (0.5*\Lx,\Ly);
        \coordinate (topLm) at (-0.5*\Lx,\Ly);
        
        \coordinate (botC) at (0,-\Ly);
        \coordinate (botR) at (\Lx,-\Ly);
        \coordinate (botL) at (-\Lx,-\Ly);
        \coordinate (botRm) at (0.5*\Lx,-\Ly);
        \coordinate (botLm) at (-0.5*\Lx,-\Ly);

        \coordinate (leftC) at (-\Lx,0);
        \coordinate (leftTm) at (-\Lx,0.5*\Ly);
        \coordinate (leftBm) at (-\Lx,-0.5*\Ly);
   
        % --- FUNDO DAS FASES ---
        \fill[DBlue] (botL) -- (botC) -- (topC) -- (topR) -- cycle;
        \fill[Blue] (botC) -- (botR) -- (topR) -- (center) -- cycle;
        \fill[LBlue] (botL) -- (center) -- (topC) -- (topL) -- cycle;

        % --- GRID E BORDAS ---
        \draw[step=4, dashed, gray] (botL) grid (topR);
        \draw[black, thick] (botL) rectangle (topR);

        % --- TEXTOS DAS FASES ---
        \node[font=\bfseries, black, yshift=-1.3em, anchor=center] at (topLm)
            {Fase paramagnética};
        \node[font=\bfseries, black, yshift=1.3em, anchor=center] at (botRm)
            {Fase ferromagnética};
        \node[font=\bfseries, white, yshift=-1.3em, anchor=center] at (topRm)
            {Fase mista};
        \node[font=\bfseries, white, yshift=1.3em, anchor=center] at (botLm)
            {Fase mista};

        % --- EIXOS ---
        \node[font=\bfseries, black, anchor=center, align=center] at ($(botC)+(0,-1)$)
            {$\boldsymbol{a}$};
        \draw[black, thick, -{Stealth[length={0.7em}, width={0.5em}]}]
            ($(botRm) + (0,-1)$) -- ($(botRm) + (4,-1)$)
            node[anchor=east, align=center, xshift=-4em] {$\boldsymbol{a>0}$};
        \draw[black, thick, -{Stealth[length={0.7em}, width={0.5em}]}]
            ($(botLm) + (0,-1)$) -- ($(botLm) + (-4,-1)$)
            node[anchor=west, align=center, xshift=4em] {$\boldsymbol{a<0}$};

        \node[font=\bfseries, black, anchor=center, align=center, rotate=90] at ($(leftC)+(-1,0)$)
            {$\boldsymbol{b}$};
        \draw[black, thick, -{Stealth[length={0.7em}, width={0.5em}]}]
            ($(leftTm) + (-1,0)$) -- ($(leftTm) + (-1,4)$)
            node[anchor=east, align=center, rotate=90,xshift=-4em] {$\boldsymbol{b>0}$};
        \draw[black, thick, -{Stealth[length={0.7em}, width={0.5em}]}]
            ($(leftBm) + (-1,0)$) -- ($(leftBm) + (-1,-4)$)
            node[anchor=west, align=center, rotate=90,xshift=4em] {$\boldsymbol{b<0}$};

        % --- MODELOS FÍSICOS ---
        \draw[Red, very thick] (center) -- (topR)
            node[below, midway, rotate={atan(\Ly/\Lx)}, font=\bfseries] {DP};
        \draw[Red, very thick] (center) -- (topC)
            node[above, midway, rotate=90, font=\bfseries] {Ising (Model A)};
        \draw[Red, very thick] (center) -- (botC)
            node[below, midway, rotate=90, font=\bfseries] {GV};
        \fill[Red] (center) circle (0.3) node[below right, font=\bfseries] {Voter Model};
        
\end{tikzpicture}

\begin{tikzpicture}[x=1em, y=1em,
        arrow/.style={-{Stealth[length=2mm, width=2mm]}, line width=1pt, black}
]

        \def\Lx{22}
        \def\Ly{12}

        \useasboundingbox (-25,-14.5) rectangle (23,12.5);

        \coordinate (center) at (0,0);

        \coordinate (topC) at (0,\Ly);
        \coordinate (topR) at (\Lx,\Ly);
        \coordinate (topL) at (-\Lx,\Ly);
        \coordinate (topRm) at (0.5*\Lx,\Ly);
        \coordinate (topLm) at (-0.5*\Lx,\Ly);
        
        \coordinate (botC) at (0,-\Ly);
        \coordinate (botR) at (\Lx,-\Ly);
        \coordinate (botL) at (-\Lx,-\Ly);
        \coordinate (botRm) at (0.5*\Lx,-\Ly);
        \coordinate (botLm) at (-0.5*\Lx,-\Ly);

        \coordinate (leftC) at (-\Lx,0);
        \coordinate (leftTm) at (-\Lx,0.5*\Ly);
        \coordinate (leftBm) at (-\Lx,-0.5*\Ly);
   
        % --- FUNDO DAS FASES ---
        \fill[DBlue] (botL) -- (botC) -- (topC) -- (topR) -- cycle;
        \fill[Blue] (botC) -- (botR) -- (topR) -- (center) -- cycle;
        \fill[LBlue] (botL) -- (center) -- (topC) -- (topL) -- cycle;

        % --- CHAMA O LUA PARA DESENHAR O FECHO ---
        \directlua{draw_convex_hull(\Lx, \Ly)}
        \node[black, font=\bfseries, rotate={atan(-\Ly/\Lx)}] at (-0.45*\Lx, 0.52*\Ly) {Domínio $q$-Voter};

        % --- GRID E BORDAS ---
        \draw[step=4, dashed, gray] (botL) grid (topR);
        \draw[black, thick] (botL) rectangle (topR);

                % --- TEXTOS DAS FASES ---
        \node[font=\bfseries, black, yshift=-1.3em, anchor=center] at (topLm)
            {Fase paramagnética};
        \node[font=\bfseries, black, yshift=1.3em, anchor=center] at (botRm)
            {Fase ferromagnética};
        \node[font=\bfseries, white, yshift=-1.3em, anchor=center] at (topRm)
            {Fase mista};
        \node[font=\bfseries, white, yshift=1.3em, anchor=center] at (botLm)
            {Fase mista};

        % --- EIXOS ---
        \node[font=\bfseries, black, anchor=center, align=center] at ($(botC)+(0,-1)$)
            {$\boldsymbol{a}$};
        \draw[black, thick, -{Stealth[length={0.7em}, width={0.5em}]}]
            ($(botRm) + (0,-1)$) -- ($(botRm) + (4,-1)$)
            node[anchor=east, align=center, xshift=-4em] {$\boldsymbol{a>0}$};
        \draw[black, thick, -{Stealth[length={0.7em}, width={0.5em}]}]
            ($(botLm) + (0,-1)$) -- ($(botLm) + (-4,-1)$)
            node[anchor=west, align=center, xshift=4em] {$\boldsymbol{a<0}$};

        \node[font=\bfseries, black, anchor=center, align=center, rotate=90] at ($(leftC)+(-1,0)$)
            {$\boldsymbol{b}$};
        \draw[black, thick, -{Stealth[length={0.7em}, width={0.5em}]}]
            ($(leftTm) + (-1,0)$) -- ($(leftTm) + (-1,4)$)
            node[anchor=east, align=center, rotate=90,xshift=-4em] {$\boldsymbol{b>0}$};
        \draw[black, thick, -{Stealth[length={0.7em}, width={0.5em}]}]
            ($(leftBm) + (-1,0)$) -- ($(leftBm) + (-1,-4)$)
            node[anchor=west, align=center, rotate=90,xshift=4em] {$\boldsymbol{b<0}$};

        % --- MODELOS FÍSICOS ---
        \draw[Red, very thick] (center) -- (topR)
            node[below, midway, rotate={atan(\Ly/\Lx)}, font=\bfseries] {DP};
        \draw[Red, very thick] (center) -- (topC)
            node[above, midway, rotate=90, font=\bfseries] {Ising (Model A)};
        \draw[Red, very thick] (center) -- (botC)
            node[below, midway, rotate=90, font=\bfseries] {GV};
        \fill[Red] (center) circle (0.3) node[below right, font=\bfseries] {Voter Model};
        
\end{tikzpicture}

\end{document}